\chapter{Effective descent morphisms}

\fontsize{12}{14.5}\selectfont

In this chapter we prove a monadicity theorem for the pullback functor. The standard terminology in the literature is \textit{descent morphism} for a morphism whose pullback functor is premonadic, and \textit{effective descent morphism} for a morphism whose pullback functor is monadic. We do not aim to develop descent theory; therefore, we will stick to the monadic theory nomenclature. 

This theorem is usually stated in the literature in the descent theory framework, and then by a comparison theorem it holds in the monadic theory. The reader can refer to \cite{JanelidzeSobralTholen2003} for the first approach. We aim to prove this monadicity theorem directly in monadic theory, without the need of the comparison theorem. This part is based on the work of Sobral and Tholen \cite{SobralTholen1992} and some ideas drawn from Janelidze and Tholen \cite{JanelidzeTholen1991}, but I will introduce some substantial changes to suit our context. Those changes will be discussed in detail in a footnote. Let us first state the theorem.

\begin{mainbox}{}
    \begin{theorem*}
In an exact category, if $p:E \to B$ is a regular epi, then $p^*$ is a monadic functor.
\end{theorem*}
\end{mainbox}

The reader may refer to \textit{\Cref{pullback functor}} for the preliminaries of this discussion. We recall that a right adjoint is monadic if the comparison functor $K: (\mathcal{A} \downarrow B) \to (\mathcal{A} \downarrow E)^T$ is an equivalence of categories, where $T=p^*p_!$ is the monad associated to $p$.
\[
% https://tikzcd.yichuanshen.de/#N4Igdg9gJgpgziAXAbVABwnAlgFyxMJZABgBpiBdUkANwEMAbAVxiRAAoAdTgWzpwAWAY0bAAggF8ABNygQA7mDoAnZQqkAhAJQgJpdJlz5CKMgCYqtRizZde-YaMkzOcxSrXypAUR16D2HgERGbklvTMrIgc3HyCIgzi0rIKSqrqvgB6ACq6ljBQAObwRKAAZmo8SGQgOBBIoVaRbADSuvogFRBViDV1SACM1BE20WiZAFQg1HACWGU4SAC0A-6dlYPU-b3D1lEgaAD6AITTILPzi4gra109jdtD53MLy08j+wCqOe3lG4hPbaNC6va7vPZsABiPwkFAkQA
\begin{tikzcd}
(\mathcal{A} \downarrow B) \arrow[rr, "K"] \arrow[dd, "p^*", shift left]                 &  & (\mathcal{A} \downarrow E)^T \arrow[lldd, "U^T", shift left] \\
                                                                                         &  &                                                              \\
(\mathcal{A} \downarrow E) \arrow[uu, "p_!", shift left] \arrow[rruu, "F^T", shift left] &  &                                                             
\end{tikzcd}
\]

Let $\mathcal{A}$ be a category with pullbacks and let $p: E \to B$ be a morphism in $\mathcal{A}$. Let $(C,\gamma; \xi)$ be an object in $(\mathcal{A} \downarrow E)^T$ with $\xi: T(C) \to C$ an algebra.

Consider the pullback in $\mathcal{A}$.
\[
% https://tikzcd.yichuanshen.de/#N4Igdg9gJgpgziAXAbVABwnAlgFyxMJZABgBpiBdUkANwEMAbAVxiRAFEAdTvAW3gD6AIQDCIAL6l0mXPkIoATOSq1GLNmMnTseAkSUBGFfWatEIIRKkgMOuUTJHqJ9efYSVMKAHN4RUABmAE4QvEgAzNQ4EEhKqqZsaFaBIWGIZCDRSAbOamYg3Lx0OAAWQbzAaCEAVuICCskgwaFIGVmIkfGuBZxFpeWVNXUGINQMdABGMAwACjK68iBBWN4lOI3NaTmZMYhxLvlo3N50vEUbqdlRu50HbMen52OT03N2eubLq+viFOJAA
\begin{tikzcd}
E\times_BC \arrow[rr, "\mathrm{proj}_{2,C}"] \arrow[d, "\mathrm{proj}_{1,C}"'] &  & C \arrow[d, "p\gamma"] \arrow[lld, "\gamma"'] \\
E \arrow[rr, "p"]                                                      &  & B                                            
\end{tikzcd}
\]
Thus, we have the span in $\mathcal{A}$:
\[
% https://tikzcd.yichuanshen.de/#N4Igdg9gJgpgziAXAbVABwnAlgFyxMJZARgBoAGAXVJADcBDAGwFcYkQBRAHS7wFt4AfQBCAYRABfUuky58hFACZSxanSat24qTOx4CRcirUMWbRCG1qYUAObwioAGYAnCHyRGQOCEjLqzdh4ADyxJaRBXd08aHyRlAM0LHj56HAALFz5gNDcAK0FFCRAaRnoAIxhGAAVZfQUQFyxbdJxJSgkgA
\begin{tikzcd}
  & E\times_BC \arrow[rd, "\xi"] \arrow[ld, "\mathrm{proj}_{2,C}"'] &   \\
C &                                                             & C
\end{tikzcd}
\]
Since $\xi$ is also a morphism in $(\mathcal{A} \downarrow E)$, we have that $\gamma \xi = \mathrm{proj}_{1,C}$
\[
% https://tikzcd.yichuanshen.de/#N4Igdg9gJgpgziAXAbVABwnAlgFyxMJZABgBpiBdUkANwEMAbAVxiRAFEAdTvAW3gD6AIQDCIAL6l0mXPkIoATOSq1GLNmMnTseAkQCMpfSvrNWiDhJUwoAc3hFQAMwBOEXkkMgcEJEtVmbNy2dLy8dBJSIK7uSGTevoj+puoW3OE4ABYuvMBobgBW4gL6INQMdABGMAwACjK68iAuWLaZOJHObh6I8T6e1CnmINwAHlhW4kA
\begin{tikzcd}
E\times_BC \arrow[rd, "\mathrm{proj}_{1,C}"'] \arrow[rr, "\xi"] &   & C \arrow[ld, "\gamma"] \\
                                                            & E &                       
\end{tikzcd}
\]

In the following for the $p^*$-image of a morphism $f$ in $(\mathcal{A}\downarrow B)$ we will use both notation $p^*(f)$ and $1_E\times_Cf$.

We recall that the unit $\eta$ of the monad $T$ in an object $(C,\gamma)$ in $(\mathcal{A} \downarrow E)$ is the morphism $\langle \gamma, 1_C \rangle$ induced by the universal property of the pullback of $p\gamma$ along $p$ and the counit $\epsilon$ in an object $(D,\zeta)$ in $(\mathcal{A} \downarrow B)$ is the morphism $\mathrm{proj}_{2,D}: E\times_BD \to D$ given by the pullback of $p$ along $\zeta$.

\begin{mainbox}{}
\begin{lemma}\label{xi coequalizer}
    Let $\mathcal{A}$ be a finitely complete category and let $p: E \to B$ be a morphism in $\mathcal{A}$. For any object $(C,\gamma;\xi)$ in $(\mathcal{A} \downarrow E)^T$, with $T$ the monad associated to $p$, the morphism $\xi$ is the coequalizer of $(1_E\times_B\xi,1_E\times_B\mathrm{proj}_{2,C})$.
\end{lemma}
\end{mainbox}

\begin{subbox}{}
\begin{proof}
  Let $(C,\gamma;\xi)$ be an object in $(\mathcal{A} \downarrow E)^T$. First we aim to show that $\xi (1_E\times_B\xi) = \xi (1_E\times_B\mathrm{proj}_{2,C})$. This holds because $\xi$ is an algebra for the monad $T=(p^*p_!, p^*\epsilon p_!, \eta)$. Observe that the multiplication in an object $(C,\gamma)$ is
  \[p^*\epsilon p_!(C,\gamma) = p^*(\epsilon_{(C,p\gamma)}) = p^*(\mathrm{proj}_{2,C}) = 1_E \times_B \mathrm{proj}_{2,C}\]
  Then we have the algebra property
  \[
  % https://tikzcd.yichuanshen.de/#N4Igdg9gJgpgziAXAbVABwnAlgFyxMJZABgBpiBdUkANwEMAbAVxiRAFEAdTvAW3gD6AIQAUXHln5xhAYQCUckAF9S6TLnyEUAJnJVajFm3F9BQmctUgM2PASJkAjPvrNWiDt1PTzltbc0iXWdqVyMPCyV9GCgAc3giUAAzACcIXiRdEBwIJABmUMN3EG4ADyw-EFT0pEdqHPzCtzYyipVktIzEMmzcxDqDZo9HAXYAAi9JMwnOcsrqrp6GxCyw4pHxyalhGd46HAALFN5gNDSAKyUBYF0ZJRBqBjoAIxgGAAV1Oy0QFKxYg44ZQUJRAA
\begin{tikzcd}
E\times_B(E\times_BC) \arrow[rr, "1_E \times_B \xi"] \arrow[d, "{1_E \times_B \mathrm{proj}_{2,C}}"'] &  & E\times_BC \arrow[d, "\xi"] \\
E\times_BC \arrow[rr, "\xi"]                                                                           &  & C                          
\end{tikzcd}
  \]
  \[\xi (1_E \times_B \mathrm{proj}_{2,C}) = \xi (1_E \times_B \xi)\]

  It remains to prove the universal property of the coequalizer. First we observe that $\xi$ is a split epi, in fact $\xi \eta_{(C,\gamma)} = 1_C$. Thus $\xi$ is the coequalizer of $(1_{E\times_BC}, \eta_{(C,\gamma)}\xi)$. We aim to prove that $\xi$ is also the coequalizer of $(1_E\times_B\xi,1_E\times_B\mathrm{proj}_{2,C})$. We already know that the universal property holds for $\xi$ with $(1_{E\times_BC}, \eta_{(C,\gamma)}\xi)$, so it suffices to show that there exist a morphism $t:E\times_BC \to E\times_B(E\times_BC)$ such that
  \begin{itemize}
    \item $(1_E\times_B\xi) t = \eta_{(C,\gamma)}\xi$
    \item $(1_E\times_B\mathrm{proj}_{2,C}) t = 1_{E\times_BC}$
  \end{itemize}
  Then given a morphism $f$ such that $f (1_E \times_B \mathrm{proj}_{2,C}) = f (1_E \times_B \xi)$, we have, composing with $t$,
  \[f (1_E \times_B \mathrm{proj}_{2,C}) t = f (1_E \times_B \xi) t \]
  \[\implies f 1_{E\times_BC} = f \eta_{(C,\gamma)}\xi \]
  By the universal property for $\xi$ with $(1_{E\times_BC}, \eta_{(C,\gamma)}\xi)$ we obtain the universal property for $\xi$ with $(1_E\times_B\xi,1_E\times_B\mathrm{proj}_{2,C})$.
  \[
  % https://tikzcd.yichuanshen.de/#N4Igdg9gJgpgziAXAbVABwnAlgFyxMJZABgBoAmAXVJADcBDAGwFcYkQBRAHS7wFt4AfQBCACm68sAuCIDCAShABfUuky58hFOQrU6TVuwn8hw2ctUgM2PASI7iehizaJOPEzLMW1NzUQBmXRpnQzcAMR8rdVstZB0ApwNXEHMlPRgoAHN4IlAAMwAnCD4kAEYaHAgkIP0XdnyoopKkMhAq8po4AAssfJwkAFoKurCQMsFjKVMeAA8sJuLSxDaOxBGevoH1kOT2CanpER4+ehxuwr5gNGKAKyVBYB1ZJRAaRnoAIxhGAAUY-xuQpYLLdAYqApLJA6drVHYgTb9IYjUIpHgwHD0R6iWSkHhZeh8U7yJRzBYQkDNZYwtYbXpI+Go-aPQ6mF5vEAfb5-AF2IEgsGLFqIGlwtpMtwDd5fH7-Px8kDA0HgyxUzqwpAAFl29TcZKFy21GsQtUYWDAKSg9B6mQ5EpA6PmcBwcAAhAACLIcrmy3laRUC8GUJRAA
\begin{tikzcd}
                                                                                                                            &  & E\times_BC \arrow[dd, "{\eta_{(C,\gamma)}\xi}", shift left] \arrow[dd, "1_{E\times_BC}"', shift right] \arrow[lldd, "t"'] &   \\
                                                                                                                            &  &                                                                                                                           &   \\
E\times_B(E\times_BC) \arrow[rr, "1_E\times_B\xi", shift left] \arrow[rr, "{1_E\times_B\mathrm{proj}_{2,C}}"', shift right] &  & E\times_BC \arrow[r, "f"] \arrow[d, "\xi"]                                                                                & F \\
                                                                                                                            &  & C \arrow[ru, "\exists! g"', dashed]                                                                                       &  
\end{tikzcd}
  \]

  We now have to prove that such a morphism $t$ actually exists. Define $t = \langle \mathrm{proj}_{1,C}, 1_{E\times_BC} \rangle$. We aim to prove that this morphism satisfies the two conditions above, ie the following diagram of pairs commutes.
  \[
  % https://tikzcd.yichuanshen.de/#N4Igdg9gJgpgziAXAbVABwnAlgFyxMJZABgBoAmAXVJADcBDAGwFcYkQBRAHS7wFt4AfQBCACm68sAuCIDCAShABfUuky58hFOQrU6TVuwn8hw2ctUgM2PASI7iehizaJOPEzLMW1NzfdIAZicDV04fK3VbLWRA3RpnQzdzJT0YKABzeCJQADMAJwg+JABGGhwIJDj9F3YePnocAAt8vmA0QoArJUFgMtklCIKi0vLKxAAWBNC6rgbm1vaunuAdAaHC4sQyEArRkDgmrFycJABaMpqkkBLBYylTHgAPLA2R7bH9w+PTxEvEsK3e7SET1RotNodCDdXprQY0Rj0ABGMEYAAUov43PksBkmqcVHlNkgdLtxpdvidzv8Zm4eDAcPReqJZKQeBl6HwGvIlM9XoSQMMtqS9n8aJTfjTam5bsBgaZ1gjkaiMX47Njcfi3sLPh8rmEeIiwBlGDAAARghaQ5a9foqM2y+VeAYWrj5ejG00gJUo9GY9UgHF4gmUJRAA
\begin{tikzcd}
                                                                                                                            &  & E\times_BC \arrow[dd, "{\eta_{(C,\gamma)}\xi}", shift left] \arrow[dd, "1_{E\times_BC}"', shift right] \arrow[lldd, "{\langle \mathrm{proj}_{1,C}, 1_{E\times_BC} \rangle}"'] &   \\
                                                                                                                            &  &                                                                                                                                                                               &   \\
E\times_B(E\times_BC) \arrow[rr, "1_E\times_B\xi", shift left] \arrow[rr, "{1_E\times_B\mathrm{proj}_{2,C}}"', shift right] &  & E\times_BC 
\end{tikzcd}
  \]
  In order to prove that we use the jointly monic pair $(\mathrm{proj}_{1,C},\mathrm{proj}_{2,C})$.
  The following compositions hold.
  \begin{itemize}
    \item $\mathrm{proj}_{2,C} (1_E\times_B\xi) \langle \mathrm{proj}_{1,C}, 1_{E\times_BC} \rangle = \xi \mathrm{proj}_{2,E\times_BC} \langle \mathrm{proj}_{1,C}, 1_{E\times_BC} \rangle = \xi 1_{E\times_BC} = \xi$
    \item $\mathrm{proj}_{2,C} \eta_{(C,\gamma)}\xi = \mathrm{proj}_{2,C} \langle \gamma, 1_C \rangle \xi = 1_C \xi = \xi$
    \item $\mathrm{proj}_{2,C} (1_E\times_B\mathrm{proj}_{2,C}) \langle \mathrm{proj}_{1,C}, 1_{E\times_BC} \rangle = \mathrm{proj}_{2,C} \mathrm{proj}_{2,E\times_BC} \langle \mathrm{proj}_{1,C}, 1_{E\times_BC} \rangle = $ \\$ = \mathrm{proj}_{2,C} 1_{E\times_BC} = \mathrm{proj}_{2,C}$
    \item $\mathrm{proj}_{2,C} 1_{E\times_BC} = \mathrm{proj}_{2,C}$
    \item $\mathrm{proj}_{1,C} (1_E\times_B\xi) \langle \mathrm{proj}_{1,C}, 1_{E\times_BC} \rangle = \mathrm{proj}_{1,E\times_BC} \langle \mathrm{proj}_{1,C}, 1_{E\times_BC} \rangle = \mathrm{proj}_{1,C}$
    \item $\mathrm{proj}_{1,C} \eta_{(C,\gamma)}\xi = \mathrm{proj}_{1,C} \langle \gamma, 1_C \rangle \xi = \gamma \xi = \mathrm{proj}_{1,C}$
    \item $\mathrm{proj}_{1,C} (1_E\times_B\mathrm{proj}_{2,C}) \langle \mathrm{proj}_{1,C}, 1_{E\times_BC} \rangle = \mathrm{proj}_{1,E\times_BC} \langle \mathrm{proj}_{1,C}, 1_{E\times_BC} \rangle = \mathrm{proj}_{1,C} $
    \item $\mathrm{proj}_{1,C} 1_{E\times_BC} = \mathrm{proj}_{1,C}$.
  \end{itemize}
  Thus $t = \langle \mathrm{proj}_{1,C}, 1_{E\times_BC} \rangle$ satisfies the conditions required.
\end{proof}
\end{subbox}

\begin{mainbox}{}
\begin{lemma}\label{equiv rel to coequalize}
    Let $\mathcal{A}$ be a finitely complete category and let $p: E \to B$ be a morphism in $\mathcal{A}$. For any object $(C,\gamma;\xi)$ in $(\mathcal{A} \downarrow E)^T$, with $T$ the monad associated to $p$, the pair $(\mathrm{proj}_{2,C},\xi)$ is an equivalence relation in $\mathcal{A}$.
\end{lemma}
\end{mainbox}

\pagebreak

\footnotetext[1]{\cite{SobralTholen1992} assumes the existence of all coequalizers, which is not the case in our setting. However, this is not a problem: at this point, we can use the morphism $p\gamma$ instead of the coequalizer of $(\mathrm{proj}_{2,C},\xi)$, as it serves our purpose. Then, \textit{a posteriori}, this proposition ensures, under the assumption of an exact category, that the pair $(\mathrm{proj}_{2,C},\xi)$ has a coequalizer, which we will make use of, following the approach of the aforementioned work.}

\begin{subbox}{}
\begin{proof}
    First we show that $(\mathrm{proj}_{2,C},\xi)$ is an internal relation. We need to show that $\langle \mathrm{proj}_{2,C},\xi \rangle$ is a mono. Consider the following diagram of solid arrows\footnotemark[1]. It commutes, in fact
    \begin{itemize}
    \item $(p\gamma)\mathrm{proj}_{2,C}=1_B(p\gamma)\mathrm{proj}_{2,C}= p \mathrm{proj}_{1,C} = p\gamma\xi$
    \item $1_B(p\gamma)p_2= 1_B(p\gamma)p_1 = p\gamma  p_1$
    \end{itemize}
    Thus we can consider the morphism $h'$ induced by the universal property of the pullback from $p_2$ and $\gamma p_1$, and $h$ similarly induced from $\mathrm{proj}_{2,C}$ and $\xi$.
    \[
% https://tikzcd.yichuanshen.de/#N4Igdg9gJgpgziAXAbVABwnAlgFyxMJZABgBpiBdUkANwEMAbAVxiRAFEAdTvAW3gD6AIQDCIAL6l0mXPkIoAjKQVVajFmxHc+g0RKkgM2PASIBmZavrNWiEGMnTjcokrNX1t+-qezTKC3dqaw07IR9DGRN5ZCUAVg8bNnYIoz8YiwTgzzZw8VUYKABzeCJQADMAJwheJAAWahwIJAA2bKS7BWEIqpqkCxAm+vbQw24iul5eOh7q2sQBocQ4ka9xyenZvuXG5sQ2tQ7DLfmlQb2BkLXOaZwAC0reYDRqgCtxAQUAchOkM6WAEyrNjcW4PJ4vCDvAQAn6OEC9eZA87DQ6jNDrKYzeGIpBkFGIM5wO5Yco4JAAWjOVzYdxA1AYWDAXigdGJhV+hN2eOoxNJ5MQVOBdjuPwZTJZbLuHJxcx5BORACMYGAoP18TS7KC6PdHs83h8AZz8UsBsrVZSzBqclrOAAPLDG7k7EDmtWCq3CkDa3UQg2fekgBh0ZUMAAKURcdkqWCKd3J+XEQA
\begin{tikzcd}
E\times_BC \arrow[rd, "h", dashed, shift left] \arrow[rrrd, "\mathrm{proj}_{2,C}", bend left] \arrow[rddd, "\xi", bend right] \arrow[rddddd, "\mathrm{proj}_{1,C}"', bend right] &                                                                                                               &  &                         \\
                                                                                                                                                                         & C\times_BC \arrow[dd, "p_1"] \arrow[rr, "p_2"] \arrow[lu, "h'", dashed, shift left] &  & C \arrow[dd, "p\gamma"] \\
                                                                                                                                                                         &                                                                                                               &  &                         \\
                                                                                                                                                                         & C \arrow[rr, "p\gamma"] \arrow[dd, "\gamma"]                                                                  &  & B \arrow[dd, "1_B"]     \\
                                                                                                                                                                         &                                                                                                               &  &                         \\
                                                                                                                                                                         & E \arrow[rr, "p"]                                                                                             &  & B                      
\end{tikzcd}
    \]
In order to prove that $\langle \mathrm{proj}_{2,C},\xi \rangle$ is a mono, we first need to show that $h$ is a mono. Since $(\mathrm{proj}_{1,C},\mathrm{proj}_{2,C})$ is a jointly monic pair, by the following compositions
\begin{itemize}
    \item $\mathrm{proj}_{1,C} h'h = \gamma p_1 h = \gamma \xi = \mathrm{proj}_{1,C} = \mathrm{proj}_{1,C} 1_{E\times_BC}$
    \item $\mathrm{proj}_{2,C} h'h = p_2 h = \mathrm{proj}_{2,C} = \mathrm{proj}_{2,C} 1_{E\times_BC}$
\end{itemize}
we have that $h'h= 1_{E\times_BC}$. This implies that $h$ is a mono, in fact suppose that $f$ and $g$ are two morphisms such that $hf=hg$, then we have
\begin{align*}
    &hf=hg \\[1.5ex]
    &\implies h'hf=h'hg \\[1.5ex]
    &\implies 1_{E\times_BC}f=1_{E\times_BC}g \\[1.5ex]
    &\implies f=g
\end{align*}

Consider now the canonical morphism $i: C\times_BC \to C\times C$ induced by the universal property of the product. We have that $i$ is a mono, in fact if $f$ and $g$ are two morphisms such that $if=ig$, then we have
\[
% https://tikzcd.yichuanshen.de/#N4Igdg9gJgpgziAXAbVABwnAlgFyxMJZARgBpiBdUkANwEMAbAVxiRAGEAdTvAW3gD6AIXYgAvqXSZc+QigDM5KrUYs2XHln5wABKIlTseAkUUAGZfWatEHcZJAYjs06QBMl1TbsHH043LIZkrUVmq2AILiyjBQAObwRKAAZgBOELxIwSA4EEhuoV5saALE9inpmYjZuUiKKtbFAm4g1Ax0AEYwDAAK-i62qVhxABY45SBpGVnUtYhkDeEgWBNTVQAss3nV1HAjWMnjiAC0C2HeyauVSJs529l7B0dnRbZxrSDtXb39JoPDYyu03mW3yhUathKxAA5ECqgs5vVzk03LCxBQxEA
\begin{tikzcd}
                                                          &                                                                   &  & C                                             \\
A \arrow[r, "f", shift left] \arrow[r, "g"', shift right] & C\times_BC \arrow[rru, "p_1"] \arrow[rrd, "p_2"'] \arrow[rr, "i"] &  & C\times C \arrow[u, "p_1'"'] \arrow[d, "p_2'"] \\
                                                          &                                                                   &  & C                                            
\end{tikzcd}
\]
\begin{itemize}
\item $if=ig \implies p_1' if = p_1' ig \implies p_1 f = p_1 g$
\item $if=ig \implies p_2' if = p_2' ig \implies p_2 f = p_2 g$
\end{itemize}
Since $(p_1,p_2)$ is a jointly monic pair, we have that $f=g$. Thus, $i$ is a mono. Since $i$ and $h$ are both monos, we have that $ih$ is a mono.

By the following commutative diagram
\[
% https://tikzcd.yichuanshen.de/#N4Igdg9gJgpgziAXAbVABwnAlgFyxMJZAJgBoBGAXVJADcBDAGwFcYkQBhAHS7wFt4AfQBCHEAF9S6TLnyEUAFgrU6TVu268sAuAAIxk6djwEiSgAwqGLNok4SpIDMblnSxK2tv3DTmSflkc2Uaa3U7AFEefiFRCRUYKABzeCJQADMAJwg+JGCQHAgkMlUbdh4+ehwAC0y+YDRsgCtxQXIAcgcM7NzEfMKkAGZQr3KuSpq6hubW4k6aRnoAIxhGAAV-VztMrCTqnC6QLJzimgHEchGyux40LDbD497hgqKLq-CQW-viEAXl1YbFymba7faPHp5M5vS6lT5YCEnRBKV5QuHeaqI3oo84vFZgKBIAC0g3yYW8FSqtXqjQgLUEv3+K3WmxBIB2ewOvieSBxbxK+MJiFJHwp4ypU1p9PI8XEQA
\begin{tikzcd}
                                                                                                                     &  &                                                                                             &  & C  \\
E\times_BC \arrow[rr, "h"] \arrow[rrrrd, "\mathrm{proj}_{2,C}"', bend right] \arrow[rrrru, "\xi", bend left] &  & C\times_BC \arrow[rru, "p_1"] \arrow[rrd, "p_2"'] \arrow[rr, "i"] &  & C\times C \arrow[d, "p_2'"'] \arrow[u, "p_1'"]           \\
                                                                                                                     &  &                                                                                             &  & C 
\end{tikzcd}
\]
we conclude that $ih = \langle \mathrm{proj}_{2,C},\xi \rangle$, so $\langle \mathrm{proj}_{2,C},\xi \rangle$ is a mono.

We proved that $(\mathrm{proj}_{2,C},\xi)$ is an internal relation, now we need to show that it is an equivalence, ie it is reflexive, symmetric and transitive.

\begin{itemize}
\item Claim: $(\mathrm{proj}_{2,C},\xi)$ is reflexive. 

Called $\eta$ the unit of the monad, we have that $\langle \mathrm{proj}_{2,C},\xi \rangle \eta_{(C,\gamma)} = \langle 1_C, 1_C \rangle$. In fact $\xi \eta_{(C,\gamma)} = 1_C$ and by the definition of the unit of the monad, we have that $\mathrm{proj}_{2,C} \eta_{(C,\gamma)} = \mathrm{proj}_{2,C} \langle \gamma, 1_C \rangle=1_C$.
\[
% https://tikzcd.yichuanshen.de/#N4Igdg9gJgpgziAXAbVABwnAlgFyxMJZARgBoAGAXVJADcBDAGwFcYkQBhEAX1PU1z5CKMgCZqdJq3YBRADpy8AW3gB9AEJde-bHgJFypcTQYs2iTjz4gMuoUVFGJp6Ra0SYUAObwioAGYAThBKSIYgOBBIZJJm7AowOPSqwAAUHKQKXvRKSvQAlNxWAcGhiOGRSI6xriDEqlw0jPQARjCMAAoCesIggVheABY4xSBBIWE0lYgAzCZS5nUNo+NlMdPVLosKeTiDgUrAaMEAVtyqoiul0VNRs-NxFgoAHlggTa3tXXb6Fv1DI24lG4QA
\begin{tikzcd}
  & C \arrow[dd, "{\eta_{(C,\gamma)}}", pos=0.7] \arrow[ldd, "1_C"'] \arrow[rdd, "1_C"] &   \\
  &                                                                            &   \\
C & E\times_BC \arrow[l, "\mathrm{proj}_{2,C}"] \arrow[r, "\xi"']                  & C
\end{tikzcd}
\]
\item Claim: $(\mathrm{proj}_{2,C},\xi)$ is symmetric. 

Consider the following diagram, where $\mathrm{proj}_{1,E\times_BC}$ and $\mathrm{proj}_{2,E\times_BC}$ are the projections of $E\times_B(E\times_BC)$. Let $l$ be the morphism induced by the universal property of the pullback from $1_{E\times_BC}$ and $\gamma\mathrm{proj}_{2,C}$, in particular it is the unique such that
\[\gamma\mathrm{proj}_{2,C}= \mathrm{proj}_{1,E\times_BC} l = \mathrm{proj}_{1,C}(1_E \times \mathrm{proj}_{2,C}) l\]
and
\[1_{E\times_BC}= \mathrm{proj}_{2,E\times_BC} l \]
\[
% https://tikzcd.yichuanshen.de/#N4Igdg9gJgpgziAXAbVABwnAlgFyxMJZABgBpiBdUkANwEMAbAVxiRAFEAdTvAW3gD6AIQDCIAL6l0mXPkIoAjKQVVajFmy48s-OMIAUWvoNEBKUxKkgM2PASIBmZavrNWiDt2N7Rl6bbkiJQcXdXdPbV1hMUl-WXtFUgBWULdNP2sZO3lkJxDqVw0PGKsbeJynFIKwtiEMsuyiMgAWVKKQGNUYKABzeCJQADMAJwheJGbqHAgkADZqtI80DJGxpCcQaaQkhfbuXjocAAth3mA0UYArcQEAJhXR8cQlTZnEW93w-cOTs4uIa53ADkQIea3eUzeOzUixA3AAHlgQNQ4EcsIMcEgALQKWIgVZPF5bRAbVHozGIHGfNgKARGHQmBFIvEEpBEt6ktEYtnUjy0+lRITfY6nc5XO7iZEgBh0ABGMAYAAUsoEPMMsD0jpiWY91pCJry4ZwDiK-lcbgopTL5UqVQlpTBuTrwdDifMYe00NwenReAcwU8Pq9tobhb8xQCbvdqNaFcqAvb1ZrMSiuRTcVZWYgyMGISB5WAoOsc4VwrTgAKTCJJc6njniS9S2wGFasGBwlA6KjugGkPW3gB2UPGn6i-6A6PSuVxu3yEBJrW9xBD3OTfMwQt94c+v10K3T20JucL7UUcRAA
\begin{tikzcd}
    \arrow[from=2-2, to=6-2, "\mathrm{proj}_{1,E\times_BC}", bend left=65, pos=0.6]
E\times_BC \arrow[rrrd, "1_{E\times_BC}", bend left] \arrow[rd, "l", dashed] \arrow[dddd, "\mathrm{proj}_{2,C}"'] &                                                                                                                                                        &  &                                                                                      \\
                                                                                                              & E\times_B(E\times_BC)) \arrow[rr, "\mathrm{proj}_{2,E\times_BC}"] \arrow[dd, "1_E\times_B\xi", shift left] \arrow[dd, "1_E\times_B\mathrm{proj}_{2,C}"', shift right] &  & E\times_BC \arrow[dd, "\xi", shift left] \arrow[dd, "\mathrm{proj}_{2,C}"', shift right] \\
                                                                                                              &                                                                                                                                                        &  &                                                                                      \\
                                                                                                              & E\times_BC \arrow[rr, "\mathrm{proj}_{2,C}", crossing over] \arrow[dd, "\mathrm{proj}_{1,C}"]                                                                                 &  & C \arrow[dd, "p\gamma"]                                                              \\
C \arrow[rd, "\gamma"']      \arrow[ur, "\eta_{(C,\gamma)}"]                                                                                 &                                                                                                                                                        &  &                                                                                      \\
                                                                                                              & E \arrow[rr, "p"]                                                                                                                                      &  & B                                                                                   
\end{tikzcd}
\]
Define $\sigma = (1_E\times_B\xi) l$. Observe that 
\[\mathrm{proj}_{2,C} \sigma = \mathrm{proj}_{2,C} (1_E\times_B\xi) l = \xi \mathrm{proj}_{2,E\times_BC} l = \xi 1_{E\times_BC} = \xi\]
It remains to show that $\xi \sigma = \mathrm{proj}_{2,C} $. First, by the following compositions and the fact that $(\mathrm{proj}_{1,C},\mathrm{proj}_{2,C})$ is a jointly monic pair
\begin{itemize}
\item $\mathrm{proj}_{1,C} (1_E\times_B\mathrm{proj}_{2,C})l = \mathrm{proj}_{1,E\times_BC} l = \gamma \mathrm{proj}_{2,C}$
\item$\mathrm{proj}_{1,C} \eta_{(C,\gamma)}\mathrm{proj}_{2,C} = \mathrm{proj}_{1,C} \langle\gamma, 1_C \rangle \mathrm{proj}_{2,C}= \gamma \mathrm{proj}_{2,C}$
\item$\mathrm{proj}_{2,C} (1_E\times_B\mathrm{proj}_{2,C})l = \mathrm{proj}_{2,C}  \mathrm{proj}_{2,E\times_BC} l = \mathrm{proj}_{2,C} 1_{E\times_BC} = \mathrm{proj}_{2,C}$
\item$\mathrm{proj}_{2,C} \eta_{(C,\gamma)}\mathrm{proj}_{2,C} = \mathrm{proj}_{2,C} \langle\gamma, 1_C \rangle \mathrm{proj}_{2,C} = 1_C \mathrm{proj}_{2,C} = \mathrm{proj}_{2,C}$
\end{itemize}
we have that $\eta_{(C,\gamma)}\mathrm{proj}_{2,C}=(1_E\times_B\mathrm{proj}_{2,C})l$.
Hence, since by \Cref{xi coequalizer} $\xi$ coequalizes $(1_E\times_B\xi,1_E\times_B\mathrm{proj}_{2,C})$, it holds that
\[\xi \sigma = \xi (1_E\times_B\xi) l = \xi (1_E\times_B\mathrm{proj}_{2,C}) l = \xi \eta_{(C,\gamma)}\mathrm{proj}_{2,C} = 1_C \mathrm{proj}_{2,C}  = \mathrm{proj}_{2,C}\]
as claimed. Thus the following diagram commutes.
\[
% https://tikzcd.yichuanshen.de/#N4Igdg9gJgpgziAXAbVABwnAlgFyxMJZAJgBoAGAXVJADcBDAGwFcYkQBRAHS7wFt4AfQBCAYRABfUuky58hFGWLU6TVu268sAuCPFSZ2PASIAWUgEYVDFm0Qh90kBiPyi5S9bV2HklTCgAc3giUAAzACcIPiQPEBwIJAsaG3V7HmxAvnoAXgAKCx5+IWEeAA8sAEpGSSdI6NiaBKQyVVt2HmycAAsIvmA0KIArCUFiWvComMQ45sQAZhTvDq4KkBpGegAjGEYABVljBRAIrEDunAmQeunk+MSFpfb0ri7e-sGIEbGrm6Smh6tVI+cpYdYgTY7faHNz2U7nS4SSgSIA
\begin{tikzcd}
  &  & E\times_BC \arrow[dd, "\sigma=(1_E\times_B\xi)l"] \arrow[rrd, "\mathrm{proj}_{2,C}"] \arrow[lld, "\xi"'] &  &   \\
C &  &                                                                                                    &  & C \\
  &  & E\times_BC \arrow[llu, "\mathrm{proj}_{2,C}"] \arrow[rru, "\xi"']                                      &  &  
\end{tikzcd}
\]
\item Claim: $(\mathrm{proj}_{2,C},\xi)$ is transitive.

By \Cref{p^*(f) is a pullback} the pullback of $\xi$ and $\mathrm{proj}_{2,C}$ is the square
\[
% https://tikzcd.yichuanshen.de/#N4Igdg9gJgpgziAXAbVABwnAlgFyxMJZABgBpiBdUkANwEMAbAVxiRAFEAdTvAW3gD6AIQAUXHln5xhAYQCUIAL6l0mXPkIoATOSq1GLNuL6ChMpSpAZseAkR0BGPfWatEIc8tU2NRMk+oXQ3djSVNPPRgoAHN4IlAAMwAnCF4kMhAcCCQHQIM3EAcBUKlhbgAPLAtElLTEXMzsxB19VzZuXjocAAsk3mA0FIArRQFgHRlFapBk1PTqLKQAZjy29w6u3v7BiBGxnRLwqa8Z2uWFppaggoqqxQpFIA
\begin{tikzcd}
E\times_B(E\times_BC) \arrow[rr, "1_E\times_B\xi"] \arrow[d, "{\mathrm{proj}_{2,E\times_BC}}"] &  & E\times_BC \arrow[d, "{\mathrm{proj}_{2,C}}"] \\
E\times_BC \arrow[rr, "\xi"]                                                                   &  & C                                            
\end{tikzcd}
\]
We aim to prove that there exists a morphism $\tau: E\times_B(E\times_BC) \to E\times_BC$ such that
\begin{itemize}
    \item $\xi (1_E\times_B\xi) = \xi \tau$
    \item $\mathrm{proj}_{2,C} \mathrm{proj}_{2,E\times_BC} = \mathrm{proj}_{2,C}\tau$
\end{itemize}
Define $\tau = (1_E\times_B\mathrm{proj}_{2,C})$, we have that
\begin{itemize}
    \item $\xi (1_E\times_B\xi) = \xi (1_E\times_B\mathrm{proj}_{2,C})$
    \item $\mathrm{proj}_{2,C} \mathrm{proj}_{2,E\times_BC} = \mathrm{proj}_{2,C} (1_E\times_B\mathrm{proj}_{2,C})$
\end{itemize}
In particular the first condition holds by \Cref{xi coequalizer}, and the second one holds by the definition of $(1_E\times_B\mathrm{proj}_{2,C})$
\end{itemize}
\end{proof}
\end{subbox}

In an exact category each equivalence relation is a kernel pair and by \Cref{equiv rel is kernel pair of the coeq} has a coequalizer. Thus in $\mathcal{A}$ exact category $(\mathrm{proj}_{2,C},\xi)$, by \Cref{equiv rel to coequalize} has a coequalizer. Let $q:C\to Q$ be such coequalizer. Observe that, considering $\mathrm{proj}_{2,C}$ and $\xi$ as morphism in $(\mathcal{A} \downarrow B)$, then $q$ is a morphism in $(\mathcal{A} \downarrow B)$ as well by the universal property. So, called $\delta$ the morphism induced by the universal property of the coequalizer from $\mathrm{proj}_{2,C}$ and $\xi$, we have that $q$ with codomain $(Q,\delta)$ is also the coequalizer of the pair in $(\mathcal{A} \downarrow B)$.
\[
% https://tikzcd.yichuanshen.de/#N4Igdg9gJgpgziAXAbVABwnAlgFyxMJZABgBpiBdUkANwEMAbAVxiRAFEAdTvAW3gD6AIQDCIAL6l0mXPkIoATKQCMVWoxZshEqSAzY8BIksrV6zVohBjJ0g3KIBmcmvOarARQlqYUAObwRKAAZgBOELxIZCA4EEjKZhqWety8dDgAFqG8wGjhAFbiAsog1Ax0AEYwDAAKMobyIKFYfhk4OiHhkYhKMXGICeoWbGjcfnS8aR0gYRFR1LFIvXAZWMHtiAC0g27J3AAeWNOz3dGLPdQraxs7SWyp6Vk5eRCFAgqlIOVVtfUOVs1Wu1bDMukhnH14mUsGBklA6CtfJ9dvdOLAGDg6McwRdIYgISirABHbziIA
\begin{tikzcd}
E\times_BC \arrow[rrd, "p\mathrm{proj}_{1,C}"'] \arrow[rr, "\xi", shift left] \arrow[rr, "\mathrm{proj}_{2,C}"', shift right] &  & C \arrow[d, "p\gamma"] \arrow[r, "q"] & Q \arrow[ld, "\delta", dashed] \\
                                                                                                                      &  & B                                     &                               
\end{tikzcd}
\]

Consider the comparison functor $K:(\mathcal{A} \downarrow B) \to (\mathcal{A} \downarrow E)^T$ defined as $K(D,\zeta)=(p^*(D,\zeta),p^*(\epsilon_{(D,\zeta}))$, where $\epsilon$ is the counit of the adjunction, and $\epsilon_{(D,\zeta)}$ is defined as the projection from $E\times_BD$ to $D$ of the pullback of $p$ along $\zeta$.

In our context of exact category, turns out that the comparison functor $K$ has a left adjoint $L:(\mathcal{A} \downarrow E)^T \to (\mathcal{A} \downarrow B)$ defined as $L(C,\gamma;\xi)=(Q,\delta)$, where $(Q,\delta)$ is the codomain of coequalizer of the pair $(\mathrm{proj}_{2,C},\xi)$ in $(\mathcal{A} \downarrow B)$.
\[
% https://tikzcd.yichuanshen.de/#N4Igdg9gJgpgziAXAbVABwnAlgFyxMJZABgBpiBdUkANwEMAbAVxiRAAoAdTgWzpwAWAY0bAAggF8ABNygQA7mDoAnZQqkAhAJQgJpdJlz5CKMgCYqtRizZde-YaMkzOcxSrXypAUR16D2HgERGbklvTMrIgc3HyCIgzi0rIKSqrqvgB6ACq6ljBQAObwRKAAZmo8SGQgOBBIoVaRbADSINRwAlhlOEgAtACM-iAVEFWINXVIA9QRNtFomQBU7SCd3b2Ig8Oj4zO19ROz1lEgaAD6AISr6z39Q-ojlQ3UU4j7t5uDx83RAKo5XSPXbTV6HRqfe4-eYgABigJ2z0QjTeNUhW32c1OABk8hIgA
\begin{tikzcd}
(\mathcal{A} \downarrow B) \arrow[rr, "K", shift left] \arrow[dd, "p^*", shift left]     &  & (\mathcal{A} \downarrow E)^T \arrow[lldd, "U^T", shift left] \arrow[ll, "L", shift left] \\
                                                                                         &  &                                                                                          \\
(\mathcal{A} \downarrow E) \arrow[uu, "p_!", shift left] \arrow[rruu, "F^T", shift left] &  &                                                                                         
\end{tikzcd}
\]

Let us introduce the counit $\rho$ and counit $\alpha$ of the adjunction $K \dashv L$. For each object $(D,\zeta)$ in $(\mathcal{A} \downarrow B)$ the counit is a morphism of the kind $\rho_{(D,\zeta)}:LK(D,\zeta) \to (D,\zeta)$. First we compute $K(D,\zeta)$.
\[K(D,\zeta) = (E\times_BD, \mathrm{proj}_{1,D}; p^*(\mathrm{proj}_{2,D}))\]
\[
% https://tikzcd.yichuanshen.de/#N4Igdg9gJgpgziAXAbVABwnAlgFyxMJZABgBpiBdUkANwEMAbAVxiRAFEAdTvAW3gD6AIQAiIAL6l0mXPkIoATOSq1GLNmMnTseAkTIBGFfWatEHCVJAYdcokqPUT680IkqYUAObwioAGYAThC8SGQgOBBISqqmbNxoWAIGlgHBoYjhkUgGTmpmIAlJCqkgQSHR1NmIAMx5ceZopeUZuRFRtfUuhZwAXjA4dO7iQA
\begin{tikzcd}
E\times_BD \arrow[d, "\mathrm{proj}_{1,D}"] \arrow[rr, "\mathrm{proj}_{2,D}"] &  & D \arrow[d, "\zeta"] \\
E \arrow[rr, "p"]                                 &  & B                   
\end{tikzcd}
\]
Then, call $q':E\times_BD \to Q'$ the coequalizer of the equivalence relation $(\mathrm{proj}_{2,E\times_BD},p^*(\mathrm{proj}_{2,D}))$, as done above we make it a coequalizer in $(\mathcal{A} \downarrow B)$ via the universal property. We obtain $q':(E\times_BD,p\mathrm{proj}_{1,D}) \to (Q',\delta')$. Then compute $LK(D,\zeta)$. By definition of the pullback functor, the top left square in the diagram below commutes. Thus we can define, as shown in the diagram, the morphism $\rho_{(D,\zeta)}$ by the universal property of the coequalizer $q'$.
\begin{equation}\label{counit rho}
% https://tikzcd.yichuanshen.de/#N4Igdg9gJgpgziAXAbVABwnAlgFyxMJZABgBoAmAXVJADcBDAGwFcYkQBRAHS7wFt4AfQBCAERABfUuky58hFOQrU6TVu3FSZ2PASJkALCoYs2iTpOkgMO+USVGaJ9eeGXtcvYtLFjas5w8-EJi7tayugokPn6m7Ny8WAJwIgAUCcEpYgCUYTaeUQYxTv7sAIqSKjBQAObwRKAAZgBOEHxIZCA4EEhKqnHmPGhYggCMYS1tHTTdSKMlAyBDI+QTre2IfbOIAMwLLtZrU4jzXT27+wE8AF4wOPRHG0Vnc5fsy4KrWiCTGwCsM3OnWcVy4fHoOAAFs0+MA0K0AFYST4gGiMegAIxgjAAChE7OZGDBGjhHkgAS9EM84JCsCSkABaU4g9hoAB6ACpUh9yLlvr9yYCkNTafSTm9BmCIdDYfCIEiUWjMdi8bYvCBmlgapDSfz1sKhYgAGwSkAARzJxsNe36Bx4sEY91RICxYCgSB2xD1xxNlOZpUlzUhEEEwFSolINzu9GyEmd6KxuPx6s12tJaKwYACUHoNOqztd7t2XsoEiAA
\begin{tikzcd}
\arrow[from=1-3, to=5-1, "\mathrm{proj}_{1,D}", bend left=20, pos=0.6]
E\times_B(E\times_BD) \arrow[dd, "\mathrm{proj}_{2,E\times_BD}"] \arrow[rr, "p^*(\mathrm{proj}_{2,D})", shift left] \arrow[rr, "\mathrm{proj}_{2,E\times_BD}"', shift right] &  & E\times_BD \arrow[dd, "\mathrm{proj}_{2,D}", pos=0.3] \arrow[rr, "q'"] &  & Q' \arrow[lldddd, "\delta'", bend left] \arrow[lldd, "{\rho_{(D,\zeta)}}"', dashed, bend left] \\
                                                                                                                                     &  &                                                &  &                                                                                              \\
E\times_BD \arrow[dd, "\mathrm{proj}_{1,D}"] \arrow[rr, "\mathrm{proj}_{2,D}", crossing over]                                                                                   &  & D \arrow[dd, "\zeta"]                          &  &                                                                                              \\
                                                                                                                                     &  &                                                &  &                                                                                              \\
E \arrow[rr, "p"]                                                                                                                    &  & B                                              &  &                                                                                            
\end{tikzcd}
\end{equation}
Observe that $\delta'$ was also defined as the morphism induced by the universal property of the coequalizer, considering the object $E\times_BD$ as an object in $(\mathcal{A} \downarrow B)$ via $p\mathrm{proj}_{1,D}$. Thus $\rho_{(D,\zeta)}$ is a morphism in $(\mathcal{A} \downarrow B)$, ie $\zeta \rho_{(D,\zeta)} = \delta'$.

For each object $(C,\gamma;\xi)$ in $(\mathcal{A}\downarrow E)^T$ the unit is a morphism of the kind $\alpha_{(C,\gamma;\xi)}:(C,\gamma;\xi) \to KL(C,\gamma;\xi)$. We compute $KL(C,\gamma;\xi)$. Let $q:C\to Q$ be the coequalizer of the equivalence relation $(\mathrm{proj}_{2,C},\xi)$ in $\mathcal{A}$. By the universal property of the coequalizer, we have that $q$ is a morphism in $(\mathcal{A} \downarrow B)$ via $p\gamma$. Thus, we can consider the coequalizer of the equivalence relation $(\mathrm{proj}_{2,C},\xi)$ in $(\mathcal{A} \downarrow B)$, which we denote by $q:(C,p\gamma) \to (Q,\delta)$. By functoriality, we have that $p^*(q)p^*(\xi) = p^*(q)p^*(\mathrm{proj}_{2,C})$. Moreover, it holds that $\xi$ is the coequalizer of $(p^*(\xi),p^*(\mathrm{proj}_{2,C}))$ by \Cref{xi coequalizer}. Thus we can define $\alpha_{(C,\gamma;\xi)}$ as the morphism induced by the universal property of the coequalizer $q$, as done in the diagram below. One can check that $\alpha_{(C,\gamma;\xi)}$ is a morphism in $(\mathcal{A} \downarrow E)^T$.
\begin{equation}\label{unit alpha}
  % https://tikzcd.yichuanshen.de/#N4Igdg9gJgpgziAXAbVABwnAlgFyxMJZABgBpiBdUkANwEMAbAVxiRAFEAdTvAW3gD6AIQAUXHln5xhAYQCUIAL6l0mXPkIoATOSq1GLNuL6ChMpSpAZseAkTJa99Zq0QduJ6WYuqbGojqO1M6GbubKvup2KAAspEH6LmwAij5WaraayHGUwQau7hJSwqkR6X7RyDoxTvlsQmnWUVlkNXlJbuxKejBQAObwRKAAZgBOELxIAOzUOBBIAGztoVZpYxNIOiBzSADM1HAAFljDOEgAtACMywXcAB5Ya+OTiFs7iPsgRydniNeJK24vDoOEOo14wDQ4wAVooBMAdDJFCBqAw6AAjGAMAAKGX8blGWD6hzOZXWLze80QMwBt04wNB4MhMLhwGuSKeGw+sypS1pbDQ3D6dF4wM5L0+7xpITpwtFdHFex5SDi-LcAEdFYhVe8+TK2NxYAwcAqyc8kGRtlT-t9Thd-vq3GgAHoAKhE9ywCjNXMt7y2jpAQJBYIhUIgsPh12MklMHJ9Lz91oOxztfxuArdHvpIaZ4cjCNIMeKZkUClRGKxuIqmhAhOJpMs5KQ-3en1tvwddTcwcZYZZ8MRyIrmJxeOidaJJK1rap7dTvyuGZ7nAeWslVIArKisGAClA6Edeii1UHOIw0Ic6PCRDJSEKRcCANyeuTDkBo0fV5psevThNINuVoqsuZ4MqGzIRqyOjJMiAHpsBiBAYGLruuq3oUIoQA
\begin{tikzcd}
E\times_B(E\times_BC) \arrow[rr, "p^*(\xi)", shift left] \arrow[dd, "{\mathrm{proj}_{2,E\times_BC}}"] \arrow[rr, "{p^*(\mathrm{proj}_{2,C})}"', shift right] &  & E\times_BC \arrow[dd, "{\mathrm{proj}_{2,C}}"', shift right] \arrow[dd, "\xi", shift left] \arrow[rr, "p^*(q)"]   &  & E\times_BQ \arrow[dd, "{\mathrm{proj}_{2,Q}}"] \\
                                                                                                                                                                      &  &                                                                                                                   &  &                                                \\
E\times_BC \arrow[rr, "\xi", shift left] \arrow[rr, "{\mathrm{proj}_{2,C}}"', shift right] \arrow[dd, "{\mathrm{proj}_{1,C}}"]                                        &  & C \arrow[dd, "p\gamma"] \arrow[lldd, "\gamma"] \arrow[rr, "q"] \arrow[rruu, "{\alpha_{(C,\gamma;\xi)}}"', dashed] &  & Q \arrow[lldd, "\delta"]                       \\
                                                                                                                                                                      &  &                                                                                                                   &  &                                                \\
E \arrow[rr, "p"]                                                                                                                                                     &  & B                                                                                                                 &  &                                               
\end{tikzcd}
\end{equation}

One can show that $\rho$ and $\alpha$ are natural transformations, and that they satisfy the universal property of the counit / unit. This follows from the universal property of the coequalizer of the equivalence relation considered.

Our goal is to show that $p^*$ is a monadic functor, ie that the comparison functor $K$ is an equivalence of categories. We prove that $K$ is an adjoint equivalence, which is the same. In order to do so, we aim to prove that the unit $\alpha$ and the counit $\rho$ of the adjunction $L \dashv K$ are pointwise isomorphisms.

\begin{mainbox}{}
\begin{prop}\label{counit iso}
    In an exact category if $p:E \to B$ is a regular epi, then the counit $\rho$ of the adjunction $L \dashv K$ is a pointwise isomorphism.
\end{prop}
\end{mainbox}

\begin{subbox}{}
\begin{proof}
    We aim to prove that for each object $(D,\zeta)$ in $(\mathcal{A} \downarrow B)$, the morphism $\rho_{(D,\zeta)}:LK(D,\zeta) \to (D,\zeta)$ is an isomorphism.
    Consider the diagram \eqref{counit rho} of the definition of the counit $\rho_{(D,\zeta)}$. An exact category is also regular, in particular regular epis are pullback stable. If $p$ is a regular epi, then $\mathrm{proj}_{2,D}$ is a regular epi as well. The top left square is a pullback by \Cref{p^*(f) is a pullback}, thus $(p^*(\mathrm{proj}_{2,D}), \mathrm{proj}_{2,E\times_BD})$ is the kernel pair of $\mathrm{proj}_{2,D}$. Since $\mathrm{proj}_{2,D}$ is a regular epi, then, by \Cref{coeq of the kernel pair}, it is the coequalizer of its kernel pair $(p^*(\mathrm{proj}_{2,D}), \mathrm{proj}_{2,E\times_BD})$. We have that $q'$ and $\mathrm{proj}_{2,D}$ are coequalizers of the same pair, thus the morphism $\rho_{(D,\zeta)}$ defined by the universal property of the coequalizer $q$ is actually an isomorphism.
\end{proof}
\end{subbox}

\begin{mainbox}{}
\begin{prop}\label{unit iso}
    In an exact category the unit $\alpha$ of the adjunction $L \dashv K$ is a pointwise isomorphism.
\end{prop}
\end{mainbox}

\begin{subbox}{}
\begin{proof}
    We aim to prove that for each object $(C, \gamma, \xi)$ in $(\mathcal{A} \downarrow E)^T$, the morphism $\alpha_{(C,\gamma;\xi)}:(C,\gamma;\xi) \to KL(C,\gamma;\xi)$ is an isomorphism. Consider the diagram \eqref{unit alpha}. An exact category is also regular, in particular regular epis are pullback stable. Since $q$ is a regular epi, then $p^*(q)$ is a regular epi as well. Moreover, since $q$ is the coequalizer of $(\xi, \mathrm{proj}_{2,C})$ which is a kernel pair, then by \Cref{kernel pair of the coeq} $(\xi, \mathrm{proj}_{2,C})$ is the kernel pair of $q$. Then, since $p^*$ is a right adjoint, it preserves kernel pairs, thus $(p^*(\xi), p^*(\mathrm{proj}_{2,C}))$ is the kernel pair of $p^*(q)$. We already proved that $p^*(q)$ is a regular epi, then by \Cref{coeq of the kernel pair} it is the coequalizer of its kernel pair $(p^*(\xi), p^*(\mathrm{proj}_{2,C}))$. By \Cref{xi coequalizer} also $\xi$ is the coequalizer of $(p^*(\xi), p^*(\mathrm{proj}_{2,C}))$. We have that $p^*(q)$ and $\xi$ are coequalizers of the same pair, thus the morphism $\alpha_{(C,\gamma;\xi)}$ defined by the universal property of the coequalizer $\xi$ is actually an isomorphism.
\end{proof}
\end{subbox}

Finally, we have the main result we were aiming to prove.

\begin{mainbox}{}
\begin{theo}\label{monadic p^*}
    In an exact category, if $p:E \to B$ is a regular epi, then $p^*$ is a monadic functor.
\end{theo}
\end{mainbox}

\begin{subbox}{}
\begin{proof}
    The exactness of the category implies that the comparison functor $K:(\mathcal{A} \downarrow B) \to (\mathcal{A} \downarrow E)^T$ has a left adjoint $L:(\mathcal{A} \downarrow E)^T \to (\mathcal{A} \downarrow B)$, as shown. Since $p$ is a regular epi, the unit $\alpha$ and the counit $\rho$ of the adjunction $L \dashv K$ are pointwise isomorphisms, by \Cref{counit iso} and \Cref{unit iso}. Thus, $L \dashv K$ is an adjoint equivalence. This implies that $K$ is an equivalence of categories. So $p^*: (\mathcal{A} \downarrow B) \to (\mathcal{A} \downarrow E)$ is a monadic functor.
\end{proof}
\end{subbox}